\documentclass[10pt,a4paper]{amsart}
\usepackage[margin=19mm]{geometry}
\usepackage{amsmath,amssymb,mathtools}
\usepackage[scr=rsfs,cal=euler]{mathalfa}
\usepackage{xfrac}
\usepackage[unicode,hidelinks,bookmarks=false]{hyperref}
\newcommand{\ie}{i.e.\ }
\newcommand{\cf}{cf.\ }
\newcommand{\resp}{respectively\ }
\newcommand{\Subsection}{\subsection}
\providecommand{\nobreakdash}{\nobreak\hbox{-}}
\setlength{\emergencystretch}{2em}
\multlinegap0pt
\usepackage{amssymb}
\usepackage{float}
\usepackage{xcolor}
\newtheorem{theorem}{Theorem}
\title{Minimal Dimensions of Maximal Commutative Matrix Algebras and Sharp Courter-like Bounds}
\author{Ma{\l}gorzata Nowak-K\k{e}pczyk}
\date{Source: arXiv:2605.01387v2 (CC BY 4.0)}
\begin{document}
\maketitle

\noindent\emph{Source and licence.} Minimal Dimensions of Maximal Commutative Matrix Algebras and Sharp Courter-like Bounds. Version arXiv:2605.01387v2 (CC BY 4.0), distributed by arXiv under the Creative Commons Attribution 4.0 International licence, http://creativecommons.org/licenses/by/4.0/, as recorded in the arXiv licence metadata for this exact version and shown on the abstract page https://arxiv.org/abs/2605.01387v2. Full licence notice: \texttt{licenses/arxiv-2605.01387-CC-BY-4.0.txt}.\par\smallskip

\noindent\textit{Editorial scope.} Counted unit: the article's theorem theorem (source0.tex lines 1259--1272). The article's complete original proof of it is delivered with it (source0.tex lines 1274--1343, one \texttt{proof} environment of 70 lines). No mathematics is written, repaired or re-derived: the statement and the proof are the article's own lines, in the article's own notation, and no retained line is substituted or re-flowed.  No further source-local context is needed: every cross-reference of every retained line resolves inside the two delivered spans.  Nothing was omitted from any retained span, so the package carries no omission record.  No retained line contains a citation, so the wrapper carries no bibliography and no bibliography entry.  No retained line contains a figure environment, an image-inclusion command or drawing code, so no image is delivered and the package declares no asset dependency.  Editorial formatting only, in addition: an AMS \texttt{amsart} preamble that restates the article's own declarations and macros used by the retained lines, namely the environments \texttt{theorem} on separate counters, together with the standard packages those lines need.  The retained environments print in this document as Theorem 1 in order of appearance, whereas the article prints the same items with the numbering of its own full document.  The complete native source of the article as distributed in the arXiv e-print for this exact version is bundled unchanged as \texttt{sources/199633/source0.tex}.
\begin{theorem}[Maximality of stacked brick algebras]
	For all integers \(q\ge 1\) and \(r\ge 0\), the algebra
	\[
	\mathcal A_{q,r}=\mathcal E^{\star q}\star\mathcal D^{\star r}
	\]
	is local maximal commutative. It has Loewy signature
	\[
	(2,5q+r,2)
	\]
	and dimension
	\[
	\dim \mathcal A_{q,r}=5+4q+r.
	\]
\end{theorem}

\begin{proof}
	By repeated application of the stack construction proposition,
	$\mathcal A_{q,r}$
	is a local commutative algebra with Loewy signature $(2,5q+r,2).$
	
	Since each copy of \(\mathcal E\) contributes four radical generators,
	while each copy of \(\mathcal D\) contributes one, and all stacked
	algebras share the same scalar identity and the same socle of dimension
	\(4\), we obtain
	\[
	\dim \mathcal A_{q,r}=1+4+4q+r=5+4q+r.
	\]
	
	It remains to prove maximal commutativity.
	
	Let
	\[
	X\in C(\mathcal A_{q,r}).
	\]
	With respect to the decomposition
	\[
	K^2
	\oplus
	K^5
	\oplus\cdots\oplus
	K^5
	\oplus
	K
	\oplus\cdots\oplus
	K
	\oplus
	K^2,
	\]
	commutation with the common socle forces \(X\) to preserve the Loewy
	filtration and to have equal scalar action on the outer diagonal blocks.
	
	Now consider the off-diagonal blocks between two distinct brick blocks.
	Such a block
	\[
	T:K^5\to K^5
	\]
	must satisfy
	\[
	TA_i=0,
	\qquad
	B_iT=0,
	\qquad i=1,2,3,4,
	\]
	because it must commute with all generators of the corresponding copy of
	\(\mathcal E\).
	By the rigidity lemma,
	\[
	T=0.
	\]
	Hence no nontrivial mixed blocks between brick components occur.
	
The same argument applies to mixed blocks involving \(\mathcal D\), which vanish
by the corresponding cases of Lemma~23. Therefore the centralizer splits
blockwise into the individual centralizers of the constituent bricks.
	
	Since $	C(\mathcal E)=\mathcal E$
	and $C(\mathcal D)=\mathcal D,$
	it follows that
	\[
	C(\mathcal A_{q,r})=\mathcal A_{q,r}.
	\]
	
	Thus $\mathcal A_{q,r}$
	is maximal commutative.
\end{proof}

\end{document}
